\section{Discussion}
\label{sec:discussion}

\subsection{Scientific validity requires more than the right comparison}

The controlled checker comparison supports a specific conclusion: matching the source--endpoint comparison is not sufficient to establish that an atomistic number answers a physical question. A record can have the right comparison ID and an appropriate unit while describing different atoms, a different reference frame, or an unlicensed mathematical relation. These properties are independent of whether the final supported/contradicted verdict happens to be correct. The 39 Qwen cases with correct verdicts but invalid proofs make that distinction observable; a final-answer-only score would count them as successes even though the submitted evidence could not justify the answer.

This result fits within established research on scientific information selection and evidence-supported answers, including SciCUEval and ALCE.\cite{yu2026scicueval,gao2023enabling} GeoScope's more limited addition is a finite, physically specified RMSD setting in which the selected source, measurement definition, and mathematical operation can be checked against explicit proof rules. It does not establish that other scientific domains can be reduced to the same rules, or that a checker capable of handling every physical inference has been built.

\subsection{Verification and evidence selection remain separate problems}

Full checking improves the correctness of accepted numerical answers, but it is a filter, not a solution strategy. When an answerable Qwen selection fails validation, the checker does not examine alternate records to recover a permitted proof. As a result, the stricter criterion removes the invalid acceptances but leaves valid numerical coverage at 28.36\%. The deterministic oracle's success on every provided case shows that poor model-based selection, rather than an absence of task-defined evidence, is the bottleneck for these answerable cases. That comparison does not imply that all real atomistic questions have available, complete, or automatically checkable evidence.

This distinction matters for evaluation. High precision among the accepted answers can be obtained by making rejection more selective, even while the system resolves few scientific questions. Reporting proof validity without coverage would overstate practical capability; reporting verdict accuracy alone would hide unsupported conclusions. Direct answers, selections, workflow rejections, and genuine scientific insufficiency should therefore remain separate outcomes. Similarly, success of a deterministic oracle using explicit rules does not amount to improved language-model reasoning. A future system that searches for alternate licensed evidence would be a different method, to be evaluated separately from the present frozen checker comparison.

The presentation and operation results help locate the remaining selection problem, but not its cause. Report rendering removed redundant metric labels as well as changing the representation, and the response format required strict evidence--operation pairings within a 160-token output budget. Controlled presentation changes, alternative selection strategies, or more flexible output handling could improve usability, but the current data do not quantify such improvements. Tests of symbolic robustness provide useful context for these questions without supplying a direct explanation for GeoScope's observed differences.\cite{mirzadeh2025gsm}

\subsection{Limits of the conclusions and future tests}

The conclusions are bounded by the benchmark's construction. The 1,760 variants are correlated templates derived from 55 source groups, not 1,760 independent physical experiments. Uncertainty therefore resamples whole source groups. Those groups are disjoint from 20 development groups but remain within the same OC20-Dense domain; they are not a demonstrated test of unseen chemical families, independent DFT accuracy, or independence from UMA or language-model training data. Two adsorbate formulas account for 38 of the 55 accepted groups. All five excluded source systems had formula C$_2$H$_3$O$_2$, and 16 accepted groups retained descriptive chemical-contact flags. UMA convergence and geometric consistency do not by themselves establish chemically accurate adsorption structures, chemisorption, catalytic activity, or a global energy minimum.

The original matched-checker study evaluated three compact instruction models under one fixed greedy, no-extended-thinking inference protocol and its stated response format. The later OC20 experiment separately evaluated four models, including Qwen3-32B, with its distinct strict raw-JSON scoring contract and a fence-normalized diagnostic. All 10,560 first attempts were retained, including 261 SmolLM outputs truncated at the fixed token cap. The protocol was frozen after development and geometry-only review, before final model inference, but was not externally preregistered. The ability of full checking to accept no invalid proofs is conditional on this task's enumerated proof rules and is not evidence of 100\% overall accuracy or an unrestricted general guarantee. The results do not establish how larger models, longer inference, different prompts, or other scientific quantities would perform.

A substantive next step is to improve the \emph{selection} of compatible and conclusive evidence, not merely to make validation stricter. Candidate methods could enumerate admissible records before selecting an operation, perform a constrained search over proofs, or ask a model to revise a rejected selection using an explicit explanation of the physical mismatch. These are proposals, not experiments conducted here. Any new model, intervention, or scoring change should be evaluated using fresh source groups rather than tuned against the inspected final set, while retaining proof validity, valid numerical coverage, rejection rates, and scientific insufficiency as distinct outcomes.
